% GENERATED FILE. Edit canonical lessons, then run npm run generate:books.
% canonical-source-hash: fnv1a64:f8c6f188279c2007
% canonical-lessons: BN-C417-not, BN-C417-etiem, BN-C417-belans, BN-C417-transphar-kora, BN-C417-chada

\chapter{Notes and Transfers}
\label{ch:bn-notes-and-transfers}
\hlchaptermodality{417}

\begin{chapteropening}
\textbf{By the end of this chapter:} \emph{I can take out money, check a balance, send money and give to a collection.}

Complete the last lesson of chapter 417: I can take out money, check a balance, send money and give to a collection.
\end{chapteropening}

\section[\texorpdfstring{\bn{নোট} (noṭ)}{নোট (noṭ)}]{\bn{নোট} (noṭ) — a banknote}

\label{lesson:BN-C417-not}



\begin{quote}
Before the new one: say the Bengali for to pay off, then the Bengali for to settle (a bill).
\end{quote}

\subsection*{You'll want to know: \bn{নোট}}
\textbf{\bn{নোট}} — \emph{noṭ} — \textquotedblleft{}a banknote\textquotedblright{}.

\begin{quote}
\bn{আমার} \bn{কাছে} \bn{একটা} \bn{একশো} \bn{টাকার} \bn{নোট} \bn{আছে।}
\end{quote}

\begin{quote}
\emph{āmār kāchhe ekṭā eksho ṭākār noṭ āchhe.}
\end{quote}

\begin{quote}
I have a hundred-taka note.
\end{quote}

\begin{cousinweb}
\textbf{\bn{নোট}} (\emph{noṭ}) is borrowed from English \emph{note}, short for \emph{banknote}.
\end{cousinweb}

\begin{grammarlens}[title={What you've built}]
One more word for getting around.
\end{grammarlens}

\subsection*{Your turn}
\begin{itemize}
\raggedright
  \item \emph{Say it:} \emph{noṭ}
  \item \emph{Say it:} \emph{noṭ}, once more
  \item \emph{Say it:} \emph{āmār kāchhe ekṭā eksho ṭākār noṭ āchhe.} — I have a hundred-taka note.
  \item \emph{Recall:} say the Bengali for to pay off, then the Bengali for to settle (a bill), then say \emph{noṭ} again
\end{itemize}

\begin{tcolorbox}[breakable,colback=teal!4,colframe=teal!35!black,title={Before you move on}]
What does \bn{নোট} mean? (\textquotedblleft{}A banknote\textquotedblright{}.) Say it once more.
\end{tcolorbox}
\section[\texorpdfstring{\bn{এটিএম} (eṭiem)}{এটিএম (eṭiem)}]{\bn{এটিএম} (eṭiem) — a cash machine; an ATM}

\label{lesson:BN-C417-etiem}



\begin{quote}
Before the new one: say the Bengali for to settle (a bill), then the Bengali for a banknote.
\end{quote}

\subsection*{You'll want to know: \bn{এটিএম}}
\textbf{\bn{এটিএম}} — \emph{eṭiem} — \textquotedblleft{}a cash machine; an ATM\textquotedblright{}.

\begin{quote}
\bn{স্টেশনের} \bn{পাশের} \bn{এটিএম} \bn{থেকে} \bn{টাকা} \bn{তুললাম।}
\end{quote}

\begin{quote}
\emph{sṭeshôner pāsher eṭiem theke ṭākā tullām.}
\end{quote}

\begin{quote}
I took money out of the cash machine beside the station.
\end{quote}

\bn{টাকা} \bn{তোলা}, to take money out; \bn{টাকা} \bn{জমা} \bn{দেওয়া}, to put it in.

\begin{cousinweb}
\textbf{\bn{এটিএম}} (\emph{eṭiem}) is borrowed from English \emph{ATM}, the letters of \emph{automated teller machine}, said as a word.
\end{cousinweb}

\begin{grammarlens}[title={What you've built}]
One more word for getting around.
\end{grammarlens}

\subsection*{Your turn}
\begin{itemize}
\raggedright
  \item \emph{Say it:} \emph{eṭiem}
  \item \emph{Say it:} \emph{eṭiem}, once more
  \item \emph{Say it:} \emph{sṭeshôner pāsher eṭiem theke ṭākā tullām.} — I took money out of the cash machine beside the station.
  \item \emph{Recall:} say the Bengali for to settle (a bill), then the Bengali for a banknote, then say \emph{eṭiem} again
\end{itemize}

\begin{tcolorbox}[breakable,colback=teal!4,colframe=teal!35!black,title={Before you move on}]
What does \bn{এটিএম} mean? (\textquotedblleft{}A cash machine; an ATM\textquotedblright{}.) Say it once more.
\end{tcolorbox}
\section[\texorpdfstring{\bn{ব্যালান্স} (bælāns)}{ব্যালান্স (bælāns)}]{\bn{ব্যালান্স} (bælāns) — a (bank) balance}

\label{lesson:BN-C417-belans}



\begin{quote}
Before the new one: say the Bengali for a banknote, then the Bengali for a cash machine.
\end{quote}

\subsection*{You'll want to know: \bn{ব্যালান্স}}
\textbf{\bn{ব্যালান্স}} — \emph{bælāns} — \textquotedblleft{}a (bank) balance\textquotedblright{}.

\begin{quote}
\bn{এটিএমে} \bn{আমি} \bn{আমার} \bn{ব্যালান্স} \bn{দেখলাম।}
\end{quote}

\begin{quote}
\emph{eṭieme āmi āmār bælāns dekhlām.}
\end{quote}

\begin{quote}
At the cash machine I checked my balance.
\end{quote}

\begin{cousinweb}
\textbf{\bn{ব্যালান্স}} (\emph{bælāns}) is borrowed from English \emph{balance}, the money left in an account.
\end{cousinweb}

\begin{grammarlens}[title={What you've built}]
One more word for getting around.
\end{grammarlens}

\subsection*{Your turn}
\begin{itemize}
\raggedright
  \item \emph{Say it:} \emph{bælāns}
  \item \emph{Say it:} \emph{bælāns}, once more
  \item \emph{Say it:} \emph{eṭieme āmi āmār bælāns dekhlām.} — At the cash machine I checked my balance.
  \item \emph{Recall:} say the Bengali for a banknote, then the Bengali for a cash machine, then say \emph{bælāns} again
\end{itemize}

\begin{tcolorbox}[breakable,colback=teal!4,colframe=teal!35!black,title={Before you move on}]
What does \bn{ব্যালান্স} mean? (\textquotedblleft{}A (bank) balance\textquotedblright{}.) Say it once more.
\end{tcolorbox}
\section[\texorpdfstring{\bn{ট্রান্সফার} \bn{করা} (ṭrānsphār kôrā)}{ট্রান্সফার করা (ṭrānsphār kôrā)}]{\bn{ট্রান্সফার} \bn{করা} (ṭrānsphār kôrā) — to transfer (money)}

\label{lesson:BN-C417-transphar-kora}



\begin{quote}
Before the new one: say the Bengali for a cash machine, then the Bengali for a (bank) balance.
\end{quote}

\subsection*{You'll want to know: \bn{ট্রান্সফার} \bn{করা}}
\textbf{\bn{ট্রান্সফার} \bn{করা}} — \emph{ṭrānsphār kôrā} — \textquotedblleft{}to transfer (money)\textquotedblright{}.

\begin{quote}
\bn{অরুণ} \bn{ফোনে} \bn{মামাকে} \bn{এক} \bn{হাজার} \bn{টাকা} \bn{ট্রান্সফার} \bn{করল।}
\end{quote}

\begin{quote}
\emph{ôrun phone māmāke ek hājār ṭākā ṭrānsphār kôrlô.}
\end{quote}

\begin{quote}
Arun transferred a thousand taka to his uncle by phone.
\end{quote}

Past (he/she …): \bn{তিনি} \bn{টাকা} \bn{ট্রান্সফার} \bn{করলেন}, \emph{tini ṭākā ṭrānsphār kôrlen}.

The one who gets the money takes -\bn{কে}: \bn{মামাকে}, to my uncle.

\begin{cousinweb}
\textbf{\bn{ট্রান্সফার} \bn{করা}} (\emph{ṭrānsphār kôrā}) joins \emph{transfer}, borrowed from English, to \bn{করা}, to do, the usual way Bengali makes a verb of an English word.
\end{cousinweb}

\begin{grammarlens}[title={What you've built}]
One more word for getting around.
\end{grammarlens}

\subsection*{Your turn}
\begin{itemize}
\raggedright
  \item \emph{Say it:} \emph{ṭrānsphār kôrā}
  \item \emph{Say it:} \emph{ṭrānsphār kôrā}, once more
  \item \emph{Say it:} \emph{ôrun phone māmāke ek hājār ṭākā ṭrānsphār kôrlô.} — Arun transferred a thousand taka to his uncle by phone.
  \item \emph{Recall:} say the Bengali for a cash machine, then the Bengali for a (bank) balance, then say \emph{ṭrānsphār kôrā} again
\end{itemize}

\begin{tcolorbox}[breakable,colback=teal!4,colframe=teal!35!black,title={Before you move on}]
What does \bn{ট্রান্সফার} \bn{করা} mean? (\textquotedblleft{}To transfer (money)\textquotedblright{}.) Say it once more.
\end{tcolorbox}
\section[\texorpdfstring{\bn{চাঁদা} (chā̃dā)}{চাঁদা (chā̃dā)}]{\bn{চাঁদা} (chā̃dā) — a contribution; a subscription}

\label{lesson:BN-C417-chada}



\begin{quote}
Before the new one: say the Bengali for a (bank) balance, then the Bengali for to transfer (money).
\end{quote}

\subsection*{You'll want to know: \bn{চাঁদা}}
\textbf{\bn{চাঁদা}} — \emph{chā̃dā} — \textquotedblleft{}a contribution; a subscription\textquotedblright{}.

\begin{quote}
\bn{পুজোর} \bn{জন্য} \bn{পাড়ার} \bn{ছেলেরা} \bn{চাঁদা} \bn{তুলল।}
\end{quote}

\begin{quote}
\emph{pujor jônyô pāṛār chhelerā chā̃dā tullô.}
\end{quote}

\begin{quote}
The local boys collected contributions for the Puja.
\end{quote}

\bn{চাঁদা} \bn{তোলা}, to collect contributions; \bn{চাঁদা} \bn{দেওয়া}, to give one.

\begin{cousinweb}
\textbf{\bn{চাঁদা}} (\emph{chā̃dā}) is usually traced to Persian \emph{chanda}, a word built on \emph{chand}, \textquotedblleft{}how much, some\textquotedblright{}; how it came to mean a contribution is not set out here.
\end{cousinweb}

\begin{grammarlens}[title={What you've built}]
One more word for getting around.
\end{grammarlens}

\subsection*{Your turn}
\begin{itemize}
\raggedright
  \item \emph{Say it:} \emph{chā̃dā}
  \item \emph{Say it:} \emph{chā̃dā}, once more
  \item \emph{Say it:} \emph{pujor jônyô pāṛār chhelerā chā̃dā tullô.} — The local boys collected contributions for the Puja.
  \item \emph{Recall:} say the Bengali for a (bank) balance, then the Bengali for to transfer (money), then say \emph{chā̃dā} again
\end{itemize}

\begin{tcolorbox}[breakable,colback=teal!4,colframe=teal!35!black,title={Before you move on}]
What does \bn{চাঁদা} mean? (\textquotedblleft{}A contribution; a subscription\textquotedblright{}.) Say it once more.
\end{tcolorbox}
